% ============================================================
\section{Background and Problem Formulation}
\label{sec:background}

This section recalls the fault-tolerant substrate our compiler targets,
formalises it as a graph, and states the placement problem.
Figure~\ref{fig:overview} previews the pipeline from a circuit to a
lattice-surgery schedule.

% ---- FIGURE: circuit -> architecture graph -> lattice surgery ----
\begin{figure*}[t]
    \centering
    \resizebox{0.8\textwidth}{!}{%
    \begin{tikzpicture}
      % ======================== (a) circuit ========================
      \begin{scope}
        \foreach \y/\name in {1.4/q_0,0.7/q_1,0/q_2}{
            \draw[ftgray] (0,\y)--(2,\y);
            \node[left,qlabel] at (-0.05,\y){$\name$};}
        % CNOT(q0,q1) at x=0.6
        \fill (0.6,1.4) circle (0.05);
        \draw (0.6,1.4)--(0.6,0.7);
        \draw (0.6,0.7) circle (0.12);
        \draw (0.48,0.7)--(0.72,0.7) (0.6,0.58)--(0.6,0.82);
        % T on q2 at x=0.6
        \node[draw,fill=white,inner sep=1.2pt,font=\footnotesize] at (0.6,0){$T$};
        % CNOT(q1,q2) at x=1.4
        \fill (1.4,0.7) circle (0.05);
        \draw (1.4,0.7)--(1.4,0);
        \draw (1.4,0) circle (0.12);
        \draw (1.28,0)--(1.52,0) (1.4,-0.12)--(1.4,0.12);
        \node[ftgray,font=\footnotesize] at (1,-0.7){(a) circuit $\C$};
      \end{scope}

      % map arrow
      \draw[-{Stealth[length=2mm]},thick,ftgray]
            (2.25,0.7)--(3.15,0.7) node[midway,above,font=\footnotesize,black]{map};

      % ==================== (b) architecture graph ====================
      \begin{scope}[xshift=3.6cm]
        \foreach \y in {0,1,2}\draw[ftgray!40] (0,\y)--(2,\y);
        \foreach \x in {0,1,2}\draw[ftgray!40] (\x,0)--(\x,2);
        \foreach \x in {0,1,2}\foreach \y in {0,1,2}\node[free] at (\x,\y){};
        \node[data] at (0,2){}; \node[qlabel] at (0,2){$q_0$};
        \node[data] at (0,0){}; \node[qlabel] at (0,0){$q_1$};
        \node[data] at (2,0){}; \node[qlabel] at (2,0){$q_2$};
        \node[magic] at (2,2){}; \node[mstar] at (2,2){};
        \node[ftgray,font=\footnotesize] at (1,-0.7){(b) mapping on $\A$};
      \end{scope}

      % route arrow
      \draw[-{Stealth[length=2mm]},thick,ftgray]
            (6.1,0.7)--(7.1,0.7) node[midway,above,font=\footnotesize,black]{route};

      % ==================== (c) lattice-surgery step ====================
      \begin{scope}[xshift=7.5cm]
        \foreach \y in {0,1,2}\draw[ftgray!40] (0,\y)--(2,\y);
        \foreach \x in {0,1,2}\draw[ftgray!40] (\x,0)--(\x,2);
        \foreach \x in {0,1,2}\foreach \y in {0,1,2}\node[free] at (\x,\y){};
        % corridors of the first layer: CNOT(q0,q1) and T(q2)
        \node[corridor] at (0,1){};                        % CNOT(q0,q1)
        \node[tile,fill=ftorange!22,draw=ftorange!55!black] at (2,1){}; % T route
        % patches
        \node[data] at (0,2){}; \node[qlabel] at (0,2){$q_0$};
        \node[data] at (0,0){}; \node[qlabel] at (0,0){$q_1$};
        \node[data] at (2,0){}; \node[qlabel] at (2,0){$q_2$};
        \node[magic] at (2,2){}; \node[mstar] at (2,2){};
        % merge (CNOT q0-q1) and T-route (q2 -> magic) highlights
        \draw[ftgreen!55!black,line width=1.1pt,shorten >=2.5mm,shorten <=2.5mm]
              (0,2)--(0,0);
        \draw[pin,ftorange!70!black] (2,0)--(2,2);
        \node[ftgray,font=\footnotesize] at (1,-0.7){(c) first layer};
      \end{scope}
    \end{tikzpicture}}
    \caption{From circuit to lattice surgery. A Clifford$+T$ circuit~(a) is
    \emph{mapped} onto the architecture graph $\A$~(b), assigning each logical
    qubit to a tile while magic states ($\star$) stay fixed. The \emph{router}
    then executes each layer by reserving vertex-disjoint corridors; panel~(c)
    shows the first layer, where a green merge corridor realises $\CNOT(q_0,q_1)$
    and an orange lane routes the $T$ gate on $q_2$ to a magic state. Choosing the
    most gates servable at once is a maximum vertex-disjoint paths problem on
    $\A$.}
    \label{fig:overview}
\end{figure*}


\subsection{Surface Codes and Lattice Surgery}

Physical qubits are too noisy to compute on directly, so fault-tolerant quantum
computation (FTQC) encodes each \emph{logical} qubit into many physical qubits
with a quantum error-correcting (QEC) code and operates on the encoded
information while errors are continually detected and corrected. The leading
choice is the \emph{surface code}, which stores one logical qubit in a
two-dimensional \emph{patch} of physical qubits; enlarging the patch---its code
distance---suppresses the logical error rate
exponentially~\cite{fowler2012surface}. A computation is compiled into the
Clifford$+T$ gate set, and each gate becomes a spatial operation on patches.

A logical two-qubit gate is realised by \emph{lattice surgery}: a strip of free
ancilla patches between the two operands is temporarily \emph{merged} with them to
measure a joint Pauli operator and then \emph{split}, the pair of operations
together implementing the gate~\cite{horsman2012lattice,litinski2019game}. The
non-Clifford $T$ gate cannot be done this way; it is teleported by consuming a
distilled \emph{magic state}, which must be delivered from a magic-state factory
to the operand patch through a free lane of the lattice~\cite{bravyi2005magic}.
Both kinds of operation therefore need a free \emph{corridor} of ancilla
patches---joining the two operands, or an operand and a magic state---for the
duration of the surgery, and these corridors are the resource the compiler must
schedule.

\subsection{Architecture Graph, Routing, and Safe Passage}

We model the processor as an \emph{architecture graph} $\A=(V,E)$ over a
rectangular $W\times H$ grid of patches. Each vertex $v\in V$ is a tile with
integer coordinates $(x,y)$; edges connect orthogonally adjacent tiles, the
corridors along which surgery can merge patches. A fixed subset
$V_{\mathrm m}\subset V$ hosts magic states, whose number and positions are
configuration parameters. At any instant a tile either holds a \emph{data
patch}---a placed logical qubit---or is \emph{free}, carrying an ancilla patch
that is reset each step and available for routing.

A circuit $\C$ over logical qubits $\Q_L$ is given in the Clifford$+T$ basis. From
it we extract, for each qubit $q$, its $T$-count $t(q)$ (its magic-state demand)
and, for each pair $(q,q')$, the two-qubit interaction count $c(q,q')$; we write
$\bar t,\sigma_t$ and $\bar c,\sigma_c$ for the mean and standard deviation of the
per-qubit $T$- and CNOT-counts. These summary statistics are the only circuit
information the placement uses (Section~\ref{sec:approach}), and the pair counts
define the \emph{CNOT interaction graph} whose density later separates sparse from
dense circuits.

Given a mapping, the router processes the circuit in layers. At each step it must
realise the active two-qubit gates by reserving vertex-disjoint paths in $\A$
between operand tiles, and the active $T$ gates by reserving a path from the
operand to a free magic-state tile; gates whose corridors would overlap are
deferred to a later step. Serving the largest set of gates at once is exactly a
maximum vertex-disjoint paths problem on $\A$, which is $\NP$-hard in
general~\cite{beverland2022surface,siraichi2018qubit}, so practical routers are
heuristic. We minimise the number of \emph{routing steps}---the lattice-surgery
time-depth---and secondarily track the average gate parallelism per step.

A placement is only useful if it stays routable as it grows. We say it satisfies
\emph{safe passage} if, after a qubit is committed, (i) every pair of
already-placed CNOT partners is still joined by some free path, (ii) every placed
patch with $T$ demand can still reach a magic state, and (iii) at least one
magic-state tile keeps a free neighbouring entry for the qubits with $T$ demand
not yet placed. We use two tests of this condition. The \emph{connectivity} test
checks reachability (i)--(iii) directly on the free lattice and admits any
placement that preserves it; it packs patches tightly and fits on smaller grids.
The \emph{cube} test is stricter, forbidding placement within the $3\times3$
neighbourhood of an occupied patch; it needs more area, but the extra spacing
leaves more routing options and yields consistently shorter routes in our
experiments, so we take \emph{cube} as the default and reserve \emph{connectivity}
for area-constrained instances.

\begin{definition}[Qubit Mapping]
    Given circuit $\C$ and architecture $\A$, a \emph{mapping}
    $\M : \Q_L \to V\setminus V_{\mathrm m}$ is an injective assignment of
    logical qubits to non-magic tiles.
\end{definition}

\begin{problem}[Routable Placement]
    Find a mapping $\M$ that (a) satisfies safe passage and (b) minimises the
    expected number of lattice-surgery routing steps required to execute $\C$
    on $\A$.
\end{problem}
